\section{Computation and Tuning}
\label{sec:computation}

\subsection{The design matrix and the solver}
\label{sec:solver}

The design \eqref{eq:design} is built one block at a time: the $n \times \NJ$
matrix of the wavelets of $U_{\md}$ is evaluated once per modulator, by the
cascade algorithm of \citet{Daubechies-Lagarias-1991} as implemented in
the \texttt{WaveBased} package \citep{Montoril-WaveBased-2026}, and then multiplied
rowwise by $X_{\cv}$. Since a compactly supported wavelet of level $\lev$ is
nonzero on an interval of length $(L-1)2^{-\lev}$, at most $L-1$ of the
$2^{\lev}$ columns of a level are nonzero at a given observation, so a block
has at most $n(L-1)J$ nonzero entries out of $n\NJ$ and the design is stored in
sparse form as soon as $J$ is moderate. Evaluating the basis once on a grid and
looking the values up removes the cascade from the inner loop of a
simulation.

Problem \eqref{eq:estimator} is then a group lasso with $p$ unpenalized coordinates,
% E5f (L10; D54): as duas citacoes do grpreg dentro dos grupos azuis de k = 4;
% a nota [VERIFICAR] de E5d saiu sem cinza, por ser nota de trabalho.
which group descent \citep{Breheny-Huang-2015} solves along a path of penalty levels; we use the
\texttt{grpreg} package \citep{Breheny-Miller-grpreg-2026} of R
\citep{RCoreTeam-2026}. Two points of scale deserve care, because they are
what ties the software to the theory. First, the solver standardizes
the columns and orthonormalizes each chunk before penalizing it
% L11 (D55, pergunta 49(b)): Simon & Tibshirani (2012), eq. (1.4) e Section 2.
\citep{Simon-Tibshirani-2012}, so that what
it computes is \eqref{eq:estimator} with $\norm{\btheta_{G}}_{2}$ replaced by
$\norm{Q_{G}^{1/2}\btheta_{G}}_{2}$, where $Q_{G}$ is the Gram matrix of the
columns of the chunk after centering, which is the estimator of Remark~S5.2 in
the Supplementary Material; and its objective,
$(2n)^{-1}\norm{\bY - \Amat\bc - \Bmat\btheta}_{2}^{2} +
\pen\sum_{G}w_{G}\norm{Q_{G}^{1/2}\btheta_{G}}_{2}$, is half that of
\eqref{eq:estimator}, so that the penalty level it reports is the $\pen$ of
\eqref{eq:estimator}, with no conversion. Second, the solver always fits an
intercept of its own, which carries the level of the constant covariate when
there is one. For the lasso option, solved by coordinate descent with the
\texttt{glmnet} package \citep{friedman2010regularization}, the penalty level reported for a
fit is the $\pen$ of the lasso problem and not the one of the solver, which
rescales its own by the ratio between the total number of columns and the
number of penalized ones when some coordinates carry zero penalty, and when one linear
covariate is the constant, its column has zero variance, and the solver may
either drop it or share the level with its own intercept; the fit is the same
either way, but the level $\lvl{\cv}$ of the constant covariate is only
recovered after the intercept is added back to it. For both
penalties the scale is checked by verifying the Karush-Kuhn-Tucker conditions
of the problem at every point of the path.

\subsection{Rescaling and the boundary}
\label{sec:boundary}

The theory places $\bU$ in $[0,1]^q$ by assumption. In practice each modulator
is mapped to that interval by a monotone transformation.
By default the image is $[0,1]$ itself, which is the setting of
Section~\ref{sec:theory}. Optionally, as in the
wavelet estimators of \citet*{Montoril-Chang-Vidakovic-2019}, the image is
$[\epsilon, 1-\epsilon]$ with $\epsilon$ fixed, which keeps the
observations away from the point where the periodization of
\eqref{eq:periodized} joins the two ends of the interval; the price is that the fitted
component integrates to zero over the rescaled range and not over $[0,1]$, a
shift of level and not of shape. The margin is a finite sample
device, and the reason to allow it is approximation: with the data away from
the ends, a component needs to be regular only on its support, where it can be
replaced by a periodic extension that matches at $0 \equiv 1$, built from a
bounded extension operator \citep[Chapter~3]{Triebel-1983}. This removes the
stall of Proposition~S3.1 and restores the bias of Lemma~\ref{lem:bias}, with
a Besov constant of order $\epsilon^{-(s-1/\pi)}$, an order that, by the
embedding of Besov into H\"older spaces, cannot be improved when the component
takes different values at the two ends of its support. Two features keep the
margin out of the theory. It has to be fixed, since with $\epsilon$ of the
order $2^{-J}$ the exponent guaranteed for the bias is no better than that of
periodization without a margin; and it degrades the design, since the
wavelets supported inside the margin are invisible to the data, and the
smallest eigenvalue of the Gram matrix of the basis restricted to the support
of the data vanishes as soon as $2^{J} \ge (L-1)/(2\epsilon)$, so that
Proposition~\ref{prop:gram} does not hold as stated. The alternative is the basis adapted to the
interval of \citet*{cohen1993wavelets}, which reproduces polynomials on $[0,1]$
and characterizes Besov regularity without any periodicity
\citep[Section 3.9]{Cohen-2003}, so that Lemma~\ref{lem:bias} holds with the
Besov space of the interval and the same proof. Its cost is structural: the
coarsest level must be large enough for the two boundary zones to be disjoint,
the constant is no longer a single column but a combination of the scaling
functions, and the constraint of \eqref{eq:additive} has to be imposed by
reparametrization, which leaves $pq(2^{\jzero}-1)$ unpenalized parameters
instead of none. We therefore keep the periodized basis as the default and
offer the interval basis as an option.

\subsection{Choosing the resolution and the penalty}
\label{sec:tuning}

The pair $(J,\pen)$ is chosen jointly. The default is $K$-fold cross-validation on a grid of resolutions, with the folds drawn once and reused for
every $J$, one design and one penalty path per resolution, and the usual two
readings of the curve: the minimizer of the cross-validated error and the
largest penalty within one standard error of it; the threshold
of Section~\ref{sec:threshold} is applied to the fit at the minimizer. In the
software the whole procedure, from the cross-validation over $(J,\pen)$ to the
threshold, is the single call \texttt{cv.wafc(x, u, y)}. Two cheaper rules are also
available for the lasso option, the Bayesian information criterion and its extended version, with
degrees of freedom equal to the number of nonzero coefficients plus the $p$
levels, which is the usual unbiased count for the lasso.

The theory supplies a fourth rule, since \eqref{eq:Jn} and \eqref{eq:lambda}
are explicit. It is not a practical competitor as it stands, because it needs
$\sigma$, which has to be estimated, and $\seff$, which is unknown, and
because it is a statement about order. In our experiments with the
lasso option the penalty
 of its theory is 7 to 13 times the one that minimizes the realized error,
and it is the penalty, more than the resolution, that makes the rule costly;
the resolution
 of its theory misses the best one by up to two levels, above
or below according to $\seff$.
% E5d: os numeros acima sao do lasso (E2.3, E1.6); com o bloco como padrao,
% passaram a dizer isso. Uma medida equivalente para o block lasso e da E5b.
% E5b: ancorar os numeros acima numa tabela da Secao 5 ou do suplemento
% (E2.3: lambda_n de 7 a 13 vezes o otimo; no nao homogeneo, J_n sozinho
% custa 2% e 18%, lambda_n sozinho 30% e 54%; J_n dois niveis abaixo com
% s' = 3/2, dois acima com s' = 1/4 em E1.6).
Its interest is as a reference for
what the cross-validated choice ought to approach.
% E5i (D68(a)): a frase das quatro regras saiu, porque o estudo da Secao 5 (D60)
% roda so a validacao cruzada do WAFC; o comentario E5b logo acima segue com a E5b.
% E5j (D72(a)): a oracao "worth measuring" saiu pelo mesmo motivo.
% E5j (D72(b)): o "$K$-fold" do 4.3 marca so o hifen em colR1; com "fold" tachado e
% "-fold" inserido, a linha passava 17 pt da margem na versao marcada.

A remark on the other penalties. The lasso option, with
$\norm{\btheta}_1$ in place of $\nGw{\btheta}$, is fit by \texttt{glmnet}; its
theory is in Section~S8 of the Supplementary Material, where its rates carry
an extra logarithmic factor, and the threshold of Section~\ref{sec:threshold}
applies to it unchanged. Replacing $\nGw{\btheta}$ in
\eqref{eq:estimator} by the sparse group lasso penalty of
\citet{Simon-Friedman-Hastie-Tibshirani-2013}, with one group per block
$(\cv,\md)$ and fit with the \texttt{sparsegl} package
\citep{Liang-Cohen-SolonHeinsfeld-Pestilli-McDonald-2024}, adds a second tuning parameter and removes entire components
rather than isolated chunks, which is what a reader interested in
knowing which modulators act on which coefficients would want, although for
 \eqref{eq:estimator} that question is already answered by thresholding the block
norms (Corollary~\ref{cor:threshold}). It is a
different estimator, with a different oracle inequality
\citep{Lounici-Pontil-vandeGeer-Tsybakov-2011}, and the theory above does not
cover it; the choice between the two is settled empirically.

% derivations/06-selecao-limiar.tex, adendo de E1.12; D45. A regra e a
% wafc_threshold(rule = "cv1se") de wafc/R/threshold.R, padrao do cv.wafc() (D48).

\subsection{Choosing the threshold}
\label{sec:threshold}

The threshold is chosen on the folds of Section~\ref{sec:tuning}, at the pair
$(J,\pen)$ already selected: each fold is refitted at that pair, with no new
search, thresholded at every candidate $t$ and scored on the observations it
left out. A thresholded fit is constant in $t$ between two consecutive block
norms, so the candidates are exhaustive and not a grid. The default is the
largest $t$ whose cross-validated error is within one standard error, between
folds, of the smallest, the one standard error rule of the penalty carried to
the threshold; the minimizer of the cross-validated error is available as an
option aimed at prediction. The fit is not refitted after the threshold: the
levels and the blocks kept are those of the fit, which is the estimator of
Corollaries~\ref{cor:threshold} and~\ref{cor:thrisk}. Neither rule has a
theorem here: Corollary~\ref{cor:threshold} gives the window of thresholds
that recovers the structure, and Corollary~\ref{cor:thrisk} what any threshold
costs.


%%%%%%%%%%%%%%%%%%%%%%%%%%%%%%%%%%%%%%%%%%%%%%%%%%%%%%%%%%%%%%%%%%%%%%%%%%%%%%
% E5h (k = 5): a Secao 5 sem os resultados. O desenho e o de D60, D62 e D63
% (docs/plano-projeto.md, E4.2), como o compendio wafc-studies/ o executa
% (config/study.yaml, core.yaml, arms.yaml, scale.yaml; R/data.R, R/methods.R,
% R/metrics.R; o gerador e wafc/R/dgp.R); os metodos e as grades sao os de D34,
% D41, D44 a D46, D56 e D63. O VCBART roda no compendio e nao entra no texto
% (D60(d)). [E4.3] espera o piloto (o modelo maior; a grade de n = 2000 e a de D69,
% aplicada na E5i);
% [E4.4], a producao. A frase do terceiro braco e a da E4.1b (ESTADO.md), com o
% pareamento escrito para p - 1 <= q, que cobre as duas celulas; a das
% covariaveis gaussianas e a de D62.

